\section{Постановка задачи}

Традиционно системы автономной навигации разрабатываются на основе высокоточных гироскопов и акселерометов. Бесплатформенная инерциальная навигационная система (БИНС), разрабатываемая в компании ЗАО \enquote{ПИК Прогресс} в которой используются оптический гироскоп и механический акселерометр, имеет себестоимость порядка нескольких миллионов рублей.
Поэтому была поставлена задача изучить возможность создания аналогичной системы автономной навигации на основе низкоточных микромеханических датчиков, входящие в состав платы микроконтроллера STM32, стоимость которой составляет около десяти тысяч рублей.


\section{Этапы разработки}
Для реализации поставленной задачи были выделены следующие этапы разработки:

\begin{enumerate}
	\item Разработка низкоуровневого ПО для микроконтроллера: создание программного обеспечения для сбора данных с датчиков, их первичной обработки и передачи по интерфейсу USB, а также реализации протокола для приёма команд управления от пользовательского приложения.
	
	\item Разработка алгоритмов фильтрации данных: исследование и реализация методов обработки сигналов с акселерометра и гироскопа, направленных на компенсацию шумов, дрейфа нуля и совмещение данных для точного определения пространственного положения.
	
	\item Создание графического интерфейса пользователя: проектирование и реализация десктопного приложения для управления микроконтроллером, визуализации поступающих данных в реальном времени.
	
	\item Разработка ПО для постобработки данных: создание отдельного программного модуля для углублённого анализа записанных измерений, формирования отчётов и выявления дефектов пути.
	
	\item Проектирование печатной платы: разработка схемы и трассировка платы для обеспечения стабильного питания микроконтроллера и подключения периферийных модулей таких как GPS-приёмник, датчик пройденного пути и прочее.
	
	\item Конструирование корпуса: проектирование защитного корпуса прибора, учитывающего условия эксплуатации на железной дороге, требования эргономики и охлаждения электронных компонентов.
	
	\item Комплексные испытания: проведение тестов разработанной системы для верификации точности измерений, оценки надёжности и соответствия техническому заданию.
\end{enumerate}

Подробное описание каждого этапа представлено в главе \ref{chap:main_part}.